\section{Síntesis de los resultados y consecuencias}
\label{nuclear:sintesis}
La estructura generadora conserva operaciones e historia antes de su realización numérica. En el transporte realizado, esa conservación se formula como una identidad de formas. Una lectura parcial distribuye la carga entre componente visible y complemento; la inclusión de ambos define una transformación recuperable. El refinamiento prolonga la misma identidad sin identificar profundidades diferentes ni sustituir el registro por su intensidad.

\subsection{Conservación, reconstrucción y dimensión}
La igualdad \(\mathcal U^*(b'\oplus b')\mathcal U=b\oplus b\) reúne lectura y memoria en una operación reversible. Para compresiones con registro separado, la telescopía conserva el terminal positivo y la serie completa de complementos. Para una holonomía unitaria fija, el terminal es la proyección sobre su espacio de vectores fijos; para una sucesión variable, lo determina el límite fuerte de los operadores de Gram.

La memoria dodecafásica permite recuperar el registro centrado mediante dos observaciones especificadas, con norma de inversa \(9/4\). La incidencia transporta esa recuperación al proyector dimensional. Su prolongación exterior conserva cada grado mediante los determinantes de Gram y exige incluir componentes mixtas. Así, las relaciones entre longitud, área y volumen se expresan mediante mapas definidos sobre el mismo registro, con sus invariantes respectivos.

\subsection{Acción y consecuencias físicas}
La acción, el reloj y la base energética mantienen sus tipos durante la composición. El cociente electrónico \(m_ec^2/\hbar\) se expresa mediante la frecuencia del reloj y el carácter electrónico, conservando el factor de conversión de base. La normalización de Maxwell--Dirac elimina las escalas dimensionales escogidas y deja el acoplamiento adimensional. El transporte gravitatorio conserva la escala areal \(\hbar G/c^3\); la solución homogénea de rebote mantiene sus condiciones de materia, torsión y constante cosmológica.

La relación de Noether se manifiesta cuando el observable transportado es una carga de la acción declarada. El cambio de forma añade un término de conexión al generador hamiltoniano y conserva exactamente la acción orientada. La holonomía de un retorno conserva, además, la información que distingue retorno de fase y retorno del estado completo.

\subsection{Memoria y espectro cuántico}
La identidad \eqref{nuclear:eq:espectro} determina el espectro simpléctico de la covarianza reducida por la componente antilineal del operador de memoria. Para el lazo
\[
H_1=\begin{pmatrix}2&1\\1&1\end{pmatrix},
\qquad
W_1=\frac19\begin{pmatrix}13&3\\3&10\end{pmatrix},
\]
se obtiene \(\det W_1=121/81\), valor simpléctico \(11/9\) y pureza \(9/11\). El espectro de la densidad reducida es
\[
\lambda_j=\frac9{10}\,10^{-j},\qquad j=0,1,2,\ldots,
\]
y su entropía satisface
\[
\frac{S}{k_B}=\frac{10}{9}\log10-\log9.
\]
Las cifras proceden del transporte y del estado fijados en el teorema. La equivalencia espectral con una densidad de Gibbs asigna una temperatura al Hamiltoniano elegido; esa equivalencia conserva su distinción respecto de un proceso dinámico de termalización.

\subsection{Contenido fundacional}
La articulación entre número, incidencia y dinámica reside en la conservación de un antecedente común y en las operaciones que relacionan sus realizaciones. El artículo hace explícita esa articulación mediante reconstrucciones, cotas, espectros y acciones. La generación de valores y el transporte de simetrías conservan una misma procedencia; sus consecuencias observables se calculan después con el lector y el estado que intervienen efectivamente. Este orden proporciona una formulación matemática precisa de la memoria dinámica y de sus realizaciones físicas.
